\documentclass[11pt,letterpaper]{article}
\usepackage[utf8]{inputenc}
\usepackage[T1]{fontenc}
\usepackage{lmodern}
\usepackage[margin=0.28in]{geometry}
\usepackage{hyperref}
\hypersetup{colorlinks=false, urlcolor=black, pdfborder=0 0 0}

\setlength{\parindent}{0pt}
\setlength{\parskip}{0pt}
\renewcommand{\baselinestretch}{0.82}
\setlength{\lineskip}{0pt}
\setlength{\lineskiplimit}{0pt}

\makeatletter
\renewcommand{\@listi}{%
  \leftmargin=\leftmargini
  \rightmargin=0pt
  \labelsep=4pt
  \labelwidth=16pt
  \itemindent=0pt
  \listparindent=0pt
  \topsep=0pt
  \partopsep=0pt
  \parsep=0pt
  \itemsep=-11pt
}
\let\saveditemize\itemize
\let\savedenditemize\enditemize
\renewcommand{\itemize}{\bgroup\setlength{\baselineskip}{0.65\baselineskip}\saveditemize}
\renewcommand{\enditemize}{\savedenditemize\egroup}
\makeatother

\begin{document}

{\Large \textbf{HSU, Ke-Cheng (Clio)}}\\[1pt]
(852) 6840 9753 \quad $|$ \quad \href{mailto:khsu@connect.ust.hk}{khsu@connect.ust.hk} \quad $|$ \quad \href{https://github.com/cliohsu1997}{github.com/cliohsu1997} \quad $|$ \quad \href{https://cliohsu1997.github.io/}{cliohsu1997.github.io}

\vspace{0.15pt}

{\large \textbf{Professional Summary}}\\[-0.5pt]
PhD candidate in Economics at HKUST (MPhil Economics) with strong training in empirical industrial organization and hands-on work on large proprietary datasets in Asia (China e-commerce; Japan retail POS). Collects and cleans large datasets, builds reproducible analysis pipelines, and turns results into figures and written summaries for research and presentation. Clear communicator in English who explains technical work to non-specialist audiences; collaborates in Mandarin and English. Hong Kong permanent resident. Seeking Data Scientist / Research Analyst (EFMAP) at the Bank for International Settlements, Hong Kong.

\vspace{0.15pt}

{\large \textbf{Education}}\\[-0.5pt]
\textbf{The Hong Kong University of Science and Technology}, Hong Kong\\
\textit{Doctor of Philosophy in Economics} \hfill Expected Graduation: June 2027\\
Research interests: Empirical industrial organization and digital economics.\\
\textit{Master of Philosophy in Economics} \hfill 2020--2022\\
CGA: 4.1/4.3.\\
\textit{Bachelor of Science in Mathematics and Economics, Minor in Actuarial Mathematics} \hfill 2015--2019\\
CGA: 3.9/4.3.\\
Awards: Dean's List (2016--2019); University Scholarship (730K HKD); Academic Achievement Medal (top 1\%); Postgraduate studentship (1.5M HKD); Redbird scholarship (40K HKD).

\vspace{0.15pt}

{\large \textbf{Experience}}\\[-0.5pt]
\textbf{Research Assistant}, National Taiwan University \hfill 2020\\
\vspace{-6pt}
\begin{itemize}
  \item Modeled crisis intervention as regulator--creditor negotiation; derived and simulated equilibrium predictions to support policy design.
\end{itemize}
\textbf{Teaching Assistant \& Lecturer}, Department of Economics, HKUST \hfill 2020--present\\
\vspace{-6pt}
\begin{itemize}
  \item Lecturer for MSc Math Camp (4.53/5); TA for PhD and undergraduate economics (avg 4.72/5); explain technical material clearly to non-specialist audiences.
\end{itemize}

\vspace{0.15pt}

{\large \textbf{Projects}}\\[-0.5pt]
\textbf{Inventory and Pricing under Resale Price Maintenance (RPM)} \textit{(w/ K.\ Kawaguchi, HKUST; K.\ Kawai, U.\ Tokyo)}\\
\vspace{-6pt}
\begin{itemize}
  \item Built and maintained analysis-ready panels from monthly Japanese bookstore POS data (deliveries, sales, returns; tens of billions of records across 1{,}000+ stores): linked titles via standardized ISBNs and merged local demographics and Amazon-scraped product attributes in Python.
  \item Estimated demand and supply with instrumental variables (2SLS) and Hausman-style price instruments; validated the estimator with Monte Carlo on simulated equilibria; selected controls with LASSO and panel fixed effects; implemented iterative equilibrium simulation in C++.
  \item Compared the dynamic model against the classical static model and showed the static version overstates marginal cost by \(\sim\)20\% once demand uncertainty and inventory decisions are ignored; presented findings through weekly descriptive statistics, regressions, and figures.
\end{itemize}
\vspace{-3pt}
\textbf{Return Policy Design in E-Commerce} \textit{(Job Market Paper)}\\
\vspace{-6pt}
\begin{itemize}
  \item Built and estimated a structural model of purchase, return-shipping-insurance take-up, and return decisions under monopoly pricing, using within-retailer switches between automatic provision and selective adoption of return-shipping insurance on Taobao (2017--2023; \(\sim\)62M RMB revenue; six apparel categories; \(\sim\)20--30\% return rates); scraped and OCR-digitized product attributes in Python.
  \item Developed a reproducible, version-controlled estimation pipeline on AWS and HPC/SLURM (Python, R, C++); cleaned and linked orders, prices, insurance, and returns; recovered category-specific demand, consumer fit uncertainty, and return responses via GMM; accelerated Monte Carlo and optimization in C++; validated policy rankings with bootstrap inference.
  \item Found a purchase--return efficiency trade-off: a stricter return policy curbs inefficient ex-post returns but raises prices and can exclude ex-ante efficient purchases; automatic provision expands valuable purchases in high-margin categories while selective adoption outperforms it by \(\sim\)15--30 pp of first-best welfare in thinner-margin categories.
\end{itemize}

\vspace{0.15pt}

{\large \textbf{Skills}}\\[-0.5pt]
\textbf{Languages:} Mandarin (Native), English (Fluent), Cantonese (Basic).\\
\textbf{Technical:} Python, R, Stata, C++, Excel, \LaTeX, AWS, HPC/SLURM; data collection and cleaning; MLE, GMM, minimum distance, EM algorithm, bootstrap.\\
\textbf{References:} Kei Kawai (UTokyo; Compass Lexecon), Kohei Kawaguchi (HKUST; UTokyo Econ Consulting), Pak Hung Au (HKUST).

\end{document}
